%% kapa: exam=2 round=34 year=2023 period=2 subject=appraisal_theory subject_ko=감정평가이론 seq=2 number=2 points=30
%% source: Q-Net 감정평가사 기출문제 (cst00309) artlSeq=5212168 file=raw/1008/5212168/2교시 23년도 감정평가사 2차시험_감정평가이론.pdf
%% conversion: kordoc 4.24.51 -> tools/gen2.py (auto); review=proofread-gpt-sol; flags=none
\begin{longquestion}{【문제 2】}
감정평가와 관련한 다음 물음에 답하시오. (30점) \par
\medskip\noindent 물음 1) 기준가치의 중요성에 대하여 설명하고, 택지비 목적의 감정평가서에 기재할 기준가치에 대하여 논하시오. (15점) \par
\medskip\noindent 물음 2) 감정평가사 甲은 한국감정평가사협회가 설치ㆍ운영하는 감정평가심사위원회의 심사위원으로서 택지비 목적의 감정평가서를 심사하고 있다. 감정평가서에 기재된 공시지가기준법상 그 밖의 요인 보정에 관한 내용은 다음의 표와 같으며, 甲은 심사결과 감정평가서의 보완이 필요하다고 판단하고 있다. 甲의 입장에서 공시지가기준법상 그 밖의 요인 보정에 있어 ‘지역요인 비교 내용의 적정성’에 대하여 세부 심사의견을 기술하시오. (15점) \par
\begin{framedtext}
1) 그 밖의 요인 보정치 산정 방법: 인근지역 또는 동일수급권 내 유사지역의 가치형성요인이 유사한 감정평가사례 중 적정한 비교사례를 선정하여 비교사례기준 비교표준지의 감정평가액과 비교표준지 공시지가에 시점수정을 한 가액의 비율을 기준으로 산정함 \par
2) 인근지역 또는 동일수급권 내 유사지역의 택지비 감정평가사례 \par
{\small\centering
\resizebox{\linewidth}{!}{%
\begin{tabular}{|c|c|c|c|c|c|c|c|}\hline
기호 & \makecell{소재지 및\\지번} & 용도지역 & \makecell{이용\\상황} & 도로조건 & 면적(m²) & \makecell{감정평가단가\\(원/m²)} & 기준시점 \\\hline
㉮ & \makecell{서울특별시\\A구 ㄱ동 65} & \makecell{제3종\\일반주거} & 아파트 & 광대소각 & 234,000 & 18,900,000 & 2022. 08. 20. \\\hline
㉯ & \makecell{서울특별시\\B구 ㄹ동 10} & \makecell{제3종\\일반주거} & 아파트 & 광대소각 & 150,000 & 21,000,000 & 2022. 09. 20. \\\hline
\end{tabular}}\par}
3) 비교사례의 선정: 감정평가사례 중 비교표준지(A구 ㄱ동 5)와 지리적으로 근접하고(A구와 B구는 서로 인접함), 토지이용계획 및 감정평가목적이 동일하거나 유사하여 비교가능성이 높은 기호 ㉯를 비교사례로 선정하였음 \par
4) 시점수정치의 산정: (감정평가서에 기재되어 있으나 생략함) \par
5) 지역요인의 비교 \par
{\scriptsize\centering
\renewcommand{\arraystretch}{1.5}
\resizebox{\linewidth}{!}{%
\begin{tabular}{|p{1.3cm}|p{2.7cm}|p{4.3cm}|c|c|p{2.3cm}|}\hline
\multirow{2}{*}{조건} & \multirow{2}{*}{항목} & \multirow{2}{*}{세항목} & \multicolumn{2}{c|}{격차율} & \multirow{2}{*}{비교 내용} \\\cline{4-5}
 & & & 사례 & 표준지 & \\\hline
\multirow{2}{*}{가로조건} & \multirow{2}{2.7cm}{가로의 폭, 구조 등의 상태} & 폭, 포장, 보도 & \multirow{2}{*}{1.00} & \multirow{2}{*}{1.00} & \multirow{2}{*}{유사함} \\\cline{3-3}
 & & 계통 및 연속성 & & & \\\hline
\multirow{3}{*}{접근조건} & 도심과의 거리 및 교통시설의 상태 & 인근교통시설의 편의성, 인근교통시설의 도시중심 접근성 & \multirow{3}{*}{1.00} & \multirow{3}{*}{1.20} & \multirow{3}{2.3cm}{표준지는 사례 대비 도시철도와의 거리 및 편익시설 배치 상태에서 우세함} \\\cline{2-3}
 & 상가의 배치상태 & 인근상가의 편의성, 인근상가의 품격 & & & \\\cline{2-3}
 & 공공 및 편익시설의 배치상태 & 학교, 공원, 병원, 관공서 등 & & & \\\hline
\multirow{6}{*}{환경조건} & 기상조건, 자연환경 & 일조, 온도, 조망, 지반, 지질 등 & \multirow{6}{*}{1.00} & \multirow{6}{*}{1.20} & \multirow{6}{2.3cm}{표준지는 사례 대비 조망 및 획지의 상태에서 우세함} \\\cline{2-3}
 & 사회환경 & 거주자의 직업, 학군 등 & & & \\\cline{2-3}
 & 획지의 상태 & 획지의 표준적인 면적, 획지의 정연성, 주변의 이용상황 등 & & & \\\cline{2-3}
 & 공급 및 처리시설의 상태 & 상수도, 하수도, 도시가스 등 & & & \\\cline{2-3}
 & 위험 및 혐오시설 & 변전소 등의 유무, 특별고압선 등의 통과 유무 & & & \\\cline{2-3}
 & 재해발생 위험성, 공해발생의 정도 & 홍수, 절벽붕괴, 소음, 대기오염 등 & & & \\\hline
\multirow{2}{1.3cm}{행정적 조건} & \multirow{2}{2.7cm}{행정상의 규제정도} & 용도지역, 지구, 구역 등 & \multirow{2}{*}{1.00} & \multirow{2}{*}{1.00} & \multirow{2}{*}{유사함} \\\cline{3-3}
 & & 기타 규제 & & & \\\hline
기타조건 & 기타 & 장래의 동향, 기타 & 1.00 & 1.00 & 유사함 \\\hline
\multicolumn{3}{|c|}{합 계} & 1.00 & 1.44 & \\\hline
\end{tabular}}\par}
6) 개별요인의 비교: (감정평가서에 기재되어 있으나 생략함) \par
7) 그 밖의 요인 보정치의 산정: (감정평가서에 기재되어 있으나 생략함) \par
\end{framedtext}
\end{longquestion}
